% !PRINTED 285
\weberpage{285}
Man findet so durch leichte, wenn auch etwas umständliche
Rechnung
\weq{15}{\begin{array}{rl@{\qquad}rl}
\omega\tau &{}= \tau^2\chi^2\omega\Theta^2, & \Theta\tau &{}= \tau^3\omega\Theta\\
\omega\tau^2 &{}= \tau\chi^2\Theta^3, & \Theta\tau^2 &{}= \tau\chi\omega\Theta^3\\
\omega\tau^3 &{}= \tau^5\chi\Theta^2, & \Theta\tau^3 &{}= \tau^5\chi\omega\\
\omega\tau^4 &{}= \tau^4\chi^2\Theta, & \Theta\tau^4 &{}= \tau^6\omega\Theta^2\\
\omega\tau^5 &{}= \tau^3\chi^2\omega\Theta, & \Theta\tau^5 &{}= \tau^2\Theta^3\\
\omega\tau^6 &{}= \tau^6\omega\Theta^3, & \Theta\tau^6 &{}= \tau^4\chi^2\Theta^2.
\end{array}}
Diese zwölf Relationen lassen sich aber alle aus vieren von
ihnen, die man auf mannigfaltige Art auswählen kann, als Folgerungen
 ableiten. Man kann z.\,B. für diese vier fundamentalen
Relationen die folgenden wählen:
\weq{16}{\begin{array}{c}
\omega\tau = \tau^2\chi^2\omega\Theta^2, \quad \omega\tau^3 = \tau^5\chi\Theta^2, \quad \omega\tau^4 = \tau^4\chi^2\Theta,\\
\Theta\tau = \tau^3\omega\Theta,
\end{array}}
aus denen mit Benutzung der Formeln (3), (12), (13) alle anderen
leicht folgen, z.\,B.:
\[
\omega\tau^5 = \tau^4\chi^2\Theta\tau = \tau^4\chi^2\tau^3\omega\Theta = \tau^3\chi^2\omega\Theta,
\]
und so die übrigen.

Wie wir früher gesehen haben, dass wir als erzeugende
Elemente der Octa\"edergruppe $\chi$ und $\Theta$ betrachten können, so
können wir jetzt als Erzeugende der Gruppe $L_7$ die zwei Elemente
 $\omega, \tau$ ansehen, denn es ergiebt sich aus (15):
\weq{17}{\Theta^3 = \omega\tau\omega\tau^6, \quad \chi = \tau^2\Theta\tau^3\omega.}
Fassen wir zusammen, so ergiebt sich folgendes Resultat:

\gesperrt{I. Sind vier Elemente $\tau, \chi, \omega, \Theta$ der Grade 7, 3, 2, 4
gegeben, bei deren Zusammensetzung die Relationen
 bestehen:}
\[
\begin{array}{ccc}
\omega\chi = \chi^2\omega, & \Theta\omega = \omega\Theta^3, & \Theta\chi = \chi^2\omega\Theta^2,\\
\Theta^2\chi = \chi\omega\Theta^3, & \omega\tau = \tau^2\chi^2\omega\Theta^2, & \omega\tau^3 = \tau^5\chi\Theta^2,\\
\omega\tau^4 = \tau^4\chi^2\Theta, & \Theta\tau = \tau^3\omega\Theta, & \chi\tau = \tau^4\chi,
\end{array}
\]
\gesperrt{so bilden die 168 Elemente}
\[
\sigma = \tau^\varrho \chi^\lambda \omega^u \Theta^v,
\]
\gesperrt{wenn $\varrho, \lambda, \mu, \nu$ volle Restsysteme nach den Moduln
 7, 3, 2, 4 durchlaufen, eine einfache Gruppe
168\textsuperscript{sten} Grades, und $\omega$ und $\tau$ können als erzeugende
 Elemente dieser Gruppe angesehen
werden.} % [sic] exponents in the enunciation text (mu, nu) do not match the exponents u, v used in the preceding formula for sigma
